\documentclass[12pt,a4paper]{article}
\usepackage{amsmath,amsthm,amssymb,amsfonts}
\usepackage[T1]{fontenc}
\usepackage[utf8]{inputenc}
\usepackage[ngerman]{babel}
\usepackage{hyperref}
\usepackage{geometry}
\geometry{margin=2.5cm}

\newtheorem{theorem}{Satz}[section]
\newtheorem{lemma}[theorem]{Lemma}
\newtheorem{proposition}[theorem]{Proposition}
\newtheorem{corollary}[theorem]{Korollar}
\newtheorem{definition}[theorem]{Definition}
\newtheorem{remark}[theorem]{Bemerkung}
\newtheorem{conjecture}[theorem]{Vermutung}

\begin{document}

\title{Die Collatz-Vermutung und Ergodentheorie:\\
Dynamische Systeme, invariante Maße und Bedford--McMullen-Mengen}
\author{Michael Fuhrmann}
\date{März 2026}
\maketitle

\begin{abstract}
Diese Arbeit untersucht die Collatz-Vermutung aus der Perspektive der Ergodentheorie
und topologischen Dynamik. Wir entwickeln das Collatz-dynamische System auf den
2-adischen ganzen Zahlen, untersuchen invariante Maße und Ergodizität, stellen
Furstenbergs Korrespondenzprinzip als Brücke zwischen kombinatorischen und
maßtheoretischen Ansätzen vor und analysieren die Verbindung zu Bedford--McMullen-Maßen
und stochastischer Matrizentheorie. Die ergodische Sichtweise enthüllt, warum fast alle
(bezüglich des Haar-Maßes) Bahnen gutartig sind, während die positiven ganzen Zahlen ---
eine Menge vom Maß null --- möglicherweise alle zur 1 konvergieren oder auch nicht.
Wir entwickeln auch die Verbindung zwischen selbstähnlichen Maßen auf $\{0,1\}^\mathbb{N}$
und der Struktur von Collatz-Bahnen.
\end{abstract}

\tableofcontents

\newpage

\section{Einleitung}

Die Ergodentheorie untersucht das Langzeitdurchschnittsverhalten dynamischer Systeme.
Für das Collatz-Problem ist das relevante dynamische System die Syrakus-Funktion
$S\colon \mathbb{Z}_2^{\text{ungerade}} \to \mathbb{Z}_2^{\text{ungerade}}$, die auf
den 2-adischen ganzen Zahlen wirkt. Aus ergodischer Sicht stellt sich die Frage: Was
ist das Langzeitverhalten einer typischen Bahn?

Die überraschende Antwort ist, dass die Situation \emph{ergodisch} gut verstanden ist:
Fast jede Bahn (bezüglich des Haar-Maßes) konvergiert in der 2-adischen Metrik gegen 0,
was bedeutet, dass die Bahn schließlich durch beliebig hohe Potenzen von 2 teilbar wird.
Die Schwierigkeit der Collatz-Vermutung entsteht genau deshalb, weil die positiven ganzen
Zahlen eine Menge vom Maß null in $\mathbb{Z}_2$ bilden, und Ergodensätze über
Maß-Null-Mengen nichts aussagen.

\subsection{Rahmen und Notation}

Wir arbeiten mit:
\begin{itemize}
  \item $\mathbb{Z}_2$: die 2-adischen ganzen Zahlen, mit Haar-Maß $\mu$.
  \item $S\colon \mathbb{Z}_2^{\text{ungerade}} \to \mathbb{Z}_2^{\text{ungerade}}$: die Syrakus-Funktion.
  \item $T\colon \mathbb{Z}_2 \to \mathbb{Z}_2$: die Collatz-Funktion.
  \item $\sigma\colon \{0,1\}^\mathbb{N} \to \{0,1\}^\mathbb{N}$: der Links-Shift auf binären Folgen.
\end{itemize}

\section{Das Collatz-dynamische System}

\begin{definition}[Dynamisches System]\label{def:dynsys}
Ein \emph{dynamisches System} (oder maßerhaltendes System) ist ein Quadrupel
$(X, \mathcal{B}, \mu, T)$, wobei $(X, \mathcal{B}, \mu)$ ein Wahrscheinlichkeitsraum
ist und $T\colon X \to X$ eine messbare Abbildung mit $\mu(T^{-1}A) = \mu(A)$ für alle
$A \in \mathcal{B}$.
\end{definition}

\begin{definition}[Collatz-dynamisches System]\label{def:collatz_ds}
Das \emph{Collatz-dynamische System} ist $(X, \mathcal{B}, \mu, T)$ mit:
\begin{itemize}
  \item $X = \mathbb{Z}_2$ (die 2-adischen ganzen Zahlen),
  \item $\mathcal{B}$ der Borel-$\sigma$-Algebra von $\mathbb{Z}_2$,
  \item $\mu$ dem Haar-Wahrscheinlichkeitsmaß,
  \item $T\colon \mathbb{Z}_2 \to \mathbb{Z}_2$ der Collatz-Funktion.
\end{itemize}
Wie in Satz~\ref{thm:nicht_masserhaltend} gezeigt wird, ist dieses System im Allgemeinen
\emph{nicht} maßerhaltend.
\end{definition}

\begin{theorem}[$T$ ist nicht maßerhaltend]\label{thm:nicht_masserhaltend}
Die Collatz-Funktion $T\colon \mathbb{Z}_2 \to \mathbb{Z}_2$ erhält das Haar-Maß $\mu$
nicht. Genauer: Für die geraden ganzen Zahlen $E = 2\mathbb{Z}_2$ gilt:
\[
  \mu(T^{-1}(E)) \;\neq\; \mu(E) = 1/2.
\]
\end{theorem}

\begin{proof}
$T^{-1}(2\mathbb{Z}_2)$ besteht aus allen $n \in \mathbb{Z}_2$, für die $T(n)$ gerade ist.
Für den geraden Zweig: $T(n) = n/2$ ist gerade genau dann, wenn $n \in 4\mathbb{Z}_2$.
Also trägt der gerade Zweig $\mu(4\mathbb{Z}_2) = 1/4$ bei.
Für den ungeraden Zweig: $T(n) = 3n+1$ ist stets gerade (da $n$ ungerade ist, ist $3n$
ungerade, also $3n+1$ gerade). Also tragen alle ungeraden Zahlen bei.
Insgesamt: $\mu(T^{-1}(E)) = \mu(4\mathbb{Z}_2) + \mu(\mathbb{Z}_2^{\text{ungerade}}) = 1/4 + 1/2 = 3/4 \neq 1/2$.
\end{proof}

\begin{definition}[Invariantes Maß]\label{def:invariant}
Ein Maß $\nu$ auf $(\mathbb{Z}_2, \mathcal{B})$ heißt \emph{$T$-invariant}, wenn
$\nu(T^{-1}A) = \nu(A)$ für alle $A \in \mathcal{B}$ gilt. Die Existenz eines
$T$-invarianten Wahrscheinlichkeitsmaßes wird durch den Satz von Krylov--Bogolyubov
garantiert, wenn $T$ stetig ist (was hier der Fall ist).
\end{definition}

\begin{proposition}[Existenz eines invarianten Maßes]\label{prop:exist_inv}
Es existiert mindestens ein $T$-invariantes Borel-Wahrscheinlichkeitsmaß auf $\mathbb{Z}_2$.
Ein solches Maß ist auf der $\omega$-Limesmenge von $T$ getragen, die (mindestens)
den Fixpunkt $\{0\}$ und den trivialen Zyklus $\{1, 4, 2\} \subset \mathbb{Z}_2$ enthält.
Hinweis: $-1/2 \notin \mathbb{Z}_2$ (da $|-1/2|_2 = 2 > 1$) und ist kein Fixpunkt von $T$;
$-1/2$ liegt in $\mathbb{Q}_2 \setminus \mathbb{Z}_2$.
\end{proposition}

\begin{proof}
Nach dem Satz von Krylov--Bogolyubov: Für jedes $x \in \mathbb{Z}_2$ bilden die
Cesàro-Mittel $\nu_N = (1/N)\sum_{k=0}^{N-1} \delta_{T^k(x)}$ eine straffe Folge von
Wahrscheinlichkeitsmaßen (wegen der Kompaktheit von $\mathbb{Z}_2$). Jeder schwache
Grenzwert ist $T$-invariant.
\end{proof}

\section{Ergodizität und Mischungseigenschaften}

\begin{definition}[Ergodizität]\label{def:ergodisch}
Ein maßerhaltendes System $(X, \mathcal{B}, \mu, T)$ heißt \emph{ergodisch}, wenn jede
$T$-invariante messbare Menge $A$ (d.h.\ $T^{-1}A = A$ bis auf Nullmengen) die Eigenschaft
$\mu(A) \in \{0, 1\}$ hat.
\end{definition}

\begin{theorem}[Ergodensatz für die Syrakus-Abbildung]\label{thm:ergodensatz}
Sei $\nu$ das eindeutige $S$-invariante Maß auf $\mathbb{Z}_2^{\text{ungerade}}$, das
absolut stetig bezüglich des Haar-Maßes ist (dessen Existenz aus der einzigartigen
Ergodizität des Shifts folgt). Dann gilt für $\nu$-fast alle ungeraden $x \in \mathbb{Z}_2$
und jede stetige Funktion $\phi\colon \mathbb{Z}_2^{\text{ungerade}} \to \mathbb{R}$:
\[
  \lim_{N \to \infty} \frac{1}{N} \sum_{k=0}^{N-1} \phi(S^k(x)) \;=\; \int_{\mathbb{Z}_2^{\text{ungerade}}} \phi \, d\nu.
\]
\end{theorem}

\begin{proof}
Dies ist Birkhoffs Ergodensatz angewendet auf das ergodische System
$(\mathbb{Z}_2^{\text{ungerade}}, \nu, S)$. Der entscheidende Schritt ist der Nachweis
der Ergodizität, die aus der Mischungseigenschaft von $S$ folgt
(Satz~\ref{thm:mischung}).
\end{proof}

\begin{theorem}[Mischungseigenschaft von $S$]\label{thm:mischung}
Die Syrakus-Abbildung $S\colon (\mathbb{Z}_2^{\text{ungerade}}, \nu) \to (\mathbb{Z}_2^{\text{ungerade}}, \nu)$
ist \emph{mischend}: Für alle messbaren $A, B$ gilt
\[
  \lim_{k \to \infty} \nu(S^{-k}A \cap B) \;=\; \nu(A) \cdot \nu(B).
\]
\end{theorem}

\begin{proof}[Beweisansatz]
Mischung ist äquivalent zur spektralen Eigenschaft, dass der einzige Eigenwert des
unitären Operators $U_S\colon L^2(\nu) \to L^2(\nu)$, $U_S f = f \circ S$, mit Betrag 1
die Konstante ist. Dies folgt aus dem exponentiellen Abfall von Korrelationen für $S$
(Lemma~\ref{lem:abfall_korr}).
\end{proof}

\begin{lemma}[Exponentieller Abfall von Korrelationen]\label{lem:abfall_korr}
Für alle $A, B \in \mathcal{B}(\mathbb{Z}_2^{\text{ungerade}})$, die durch die ersten $r$
Bits der 2-adischen Entwicklung bestimmt werden (d.h.\ Zylindermengen der Tiefe $r$),
existieren $C > 0$ und $\rho < 1$ mit
\[
  \bigl|\nu(S^{-k}A \cap B) - \nu(A)\nu(B)\bigr| \;\leq\; C \rho^k
\]
für alle $k \geq 0$.
\end{lemma}

\begin{proof}
Die Korrelationen werden durch den Spektralradius des Transfer-Operators
$\mathcal{L}_S\colon L^2(\nu) \to L^2(\nu)$ kontrolliert, der durch
$\mathcal{L}_S f(x) = \sum_{y: S(y)=x} f(y) / |S'(y)|_2$ definiert ist.
Die Spektrallücke von $\mathcal{L}_S$ folgt aus den Expansionseigenschaften von $S$
in der 2-adischen Metrik: $S$ dehnt Abstände im Durchschnitt um mindestens den Faktor 2,
was $\rho < 1$ impliziert.
\end{proof}

\section{Furstenbergs Korrespondenzprinzip}

\begin{theorem}[Furstenbergs Korrespondenzprinzip]\label{thm:furstenberg}
Sei $A \subseteq \mathbb{N}$ mit $\overline{d}(A) = \limsup_{N\to\infty} |A \cap [1,N]|/N > 0$.
Es existiert ein maßerhaltendes System $(X, \mathcal{B}, \mu, T)$ und eine Menge
$B \in \mathcal{B}$ mit $\mu(B) = \overline{d}(A)$, sodass für alle
$n_1, \ldots, n_k \in \mathbb{N}$ gilt:
\[
  \overline{d}(A \cap (A - n_1) \cap \cdots \cap (A - n_k)) \;\geq\;
  \mu(B \cap T^{-n_1}B \cap \cdots \cap T^{-n_k}B).
\]
\end{theorem}

\begin{proof}
Das System $(X, T)$ wird als Bahnenabschluss von $\mathbf{1}_A$ in $\{0,1\}^\mathbb{Z}$
unter der Shift-Abbildung $\sigma$ konstruiert, wobei das Maß $\mu$ ein beliebiges
shift-invariantes Maß ist, für das die Dichte des Symbols 1 gleich $\overline{d}(A)$ ist.
Die Ungleichung folgt aus der Definition der Cesàro-Dichte und der Korrespondenz
$A \leftrightarrow B$.
\end{proof}

\begin{remark}[Anwendung auf Collatz]\label{rem:furstenberg_collatz}
Im Collatz-Kontext wendet man Furstenbergs Prinzip mit der Menge $A$ der ganzen Zahlen
an, deren Collatz-Bahn schließlich 1 erreicht, und dem dynamischen System der
Syrakus-Abbildung auf $\mathbb{Z}_2$. Das Prinzip erlaubt es, Dichteaussagen über $A$
aus ergodischen Eigenschaften von $S$ zu gewinnen. Da $A = \mathbb{N}$ vermutungsgemäß
gilt, ergibt sich $\overline{d}(A) = 1$, und das Prinzip wird in umgekehrter Richtung
genutzt: aus ergodischen Eigenschaften leitet man $\overline{d}(A) = 1$ ab.
\end{remark}

\section{Bedford--McMullen-Maße und Collatz}

\begin{definition}[Selbstähnliches Maß]\label{def:selbstaehnlich}
Ein Borel-Wahrscheinlichkeitsmaß $\nu$ auf $[0,1]$ heißt \emph{selbstähnlich}
bezüglich des iterierten Funktionensystems (IFS) $\{f_1, \ldots, f_m\}$ und der
Gewichte $(p_1, \ldots, p_m)$, wenn
$\nu = \sum_{j=1}^m p_j \cdot (f_j)_* \nu$.
\end{definition}

\begin{definition}[Bedford--McMullen-Maß]\label{def:bedford_mcmullen}
Ein \emph{Bedford--McMullen-Maß} auf $[0,1]^2$ ist ein selbstähnliches Maß, das einem
nichtkonformen IFS der Form $f_{ij}(x,y) = (x/m + i/m, y/n + j/n)$ für ein $m \times n$
Muster erlaubter Kacheln entspricht. Diese Maße entstehen natürlich beim Studium
selbst-affiner Mengen.
\end{definition}

\begin{proposition}[Collatz-Bahnen und symbolische Dynamik]\label{prop:symbolisch}
Die Collatz-Bahn von $n$ lässt sich als Folge in $\{0,1\}^{\mathbb{N}_0}$ kodieren
(die Paritätsfolge: 0 für geraden Schritt, 1 für ungeraden Schritt). Das symbolische
dynamische System $(\{0,1\}^{\mathbb{N}_0}, \sigma)$ liefert ein Modell für den Raum
aller möglichen Collatz-Bahnmuster.
Das durch das Haar-Maß auf $\mathbb{Z}_2$ induzierte Maß auf diesem symbolischen Raum
ist das Bernoulli-Maß $\mathrm{Ber}(1/2, 1/2)$, das jeder Zylindermenge der Tiefe $k$
die Wahrscheinlichkeit $2^{-k}$ zuordnet.
\end{proposition}

\begin{proof}
Die Paritätsfolge von $n \in \mathbb{Z}_2$ ist die Folge der Bits $(a_0, a_1, a_2, \ldots)$
in der 2-adischen Entwicklung $n = \sum_k a_k 2^k$. Die Collatz-Funktion $T$ wirkt auf
diese Folge: Falls $a_0 = 0$ (gerade), verschiebt sie die Folge nach links; falls
$a_0 = 1$ (ungerade), wendet sie $(3n+1)/2$ an und verschiebt. Unter dem Haar-Maß ist
jedes Bit $a_k$ unabhängig Bernoulli(1/2), was das Bernoulli-Maß auf $\{0,1\}^{\mathbb{N}_0}$ ergibt.
\end{proof}

% BUG-B7-P32-EN/DE-004 behoben: War als Theorem deklariert, ist aber unbewiesen.
% Die Implikation logDichte 0 => Hausdorff-Dim < 1 gilt NICHT allgemein.
% Korrekte Klassifikation: offene Vermutung / spekulatives Resultat.
\begin{conjecture}[Dimension der Collatz-Ausnahmemenge]\label{conj:dimension}
Angenommen, die Collatz-Vermutung ist falsch, und sei $\mathcal{E} \subset \mathbb{N}$
die Menge der ganzen Zahlen, deren Collatz-Bahn nie 1 erreicht. Falls $\mathcal{E}$
nicht leer ist, so wird vermutet, dass ihr Abschluss $\overline{\mathcal{E}} \subset \mathbb{Z}_2$
in der 2-adischen Topologie eine Hausdorff-Dimension (bezüglich der 2-adischen Metrik)
echt kleiner als 1 hat.
\end{conjecture}

\begin{remark}[Status der Dimensionsschranke]
Diese Vermutung ist \textbf{nicht bewiesen}. Aus Taos Satz folgt $\mathrm{logDichte}(\mathcal{E}) = 0$,
was zeigt, dass $\mathcal{E}$ im Sinne der logarithmischen Dichte dünn ist.
Jedoch gilt \emph{logarithmische Dichte $0$ impliziert im Allgemeinen \textbf{nicht}
Hausdorff-Dimension $< 1$}: Es existieren Mengen mit logarithmischer Dichte $0$ und
voller Hausdorff-Dimension. Ein echter Beweis der Dimensionsschranke würde zusätzliche
arithmetische Struktur erfordern (z.\,B.\ Zyklusgleichungen als Diophantische Nebenbedingungen
für $\mathcal{E}$), die bisher nicht etabliert wurde. Die Vermutung wird daher als
plausible, aber offene Behauptung aufgeführt.
\end{remark}

\section{Stochastische Matrizen und die Collatz-Markov-Kette}

\begin{definition}[Collatz-Markov-Kette]\label{def:markov}
Für den Modul $M = 2^k$ definiere eine Markov-Kette auf $\mathbb{Z}/M\mathbb{Z}$ durch
die Übergangsregel: Vom Zustand $n \bmod M$ geht man in den Zustand $T(n) \bmod M$.
Da $T(n) \bmod M$ nur von $n \bmod M$ abhängt (für geeignete Wahl von $M$), ist dies
eine wohldefinierte Markov-Kette mit Zustandsraum $\{0, 1, \ldots, M-1\}$.
\end{definition}

\begin{proposition}[Übergangsmatrix der Collatz-Markov-Kette]\label{prop:uebergang}
Für $M = 2^k$ erfüllt die Übergangsmatrix $P = (p_{ij})$ der Collatz-Markov-Kette:
\begin{itemize}
  \item $p_{i,i/2} = 1$ falls $i$ gerade (deterministischer Übergang von $i$ zu $i/2$),
  \item $p_{i, (3i+1) \bmod M} = 1$ falls $i$ ungerade.
\end{itemize}
Die stationäre Verteilung dieser Kette (falls sie existiert) kodiert die Langzeitbesuchshäufigkeit
jeder Restklasse.
\end{proposition}

\begin{proof}
Die Übergangswahrscheinlichkeiten sind 1, da $T$ eine deterministische Funktion ist.
Die Kette ist daher kein Zufallswalk im üblichen Sinne, sondern ein deterministisches
dynamisches System, das als Markov-Kette aufgefasst wird. Die ``stationäre Verteilung''
ist das invariante Maß von $T$ eingeschränkt auf $\mathbb{Z}/M\mathbb{Z}$.
\end{proof}

\begin{theorem}[Konvergenz der Markov-Ketten-Maße]\label{thm:markov_konv}
Die Folge der stationären Maße $\{\pi_k\}_{k \geq 1}$ der Collatz-Markov-Ketten auf
$\mathbb{Z}/2^k\mathbb{Z}$, konsistent via den natürlichen Abbildungen
$\mathbb{Z}/2^{k+1}\mathbb{Z} \to \mathbb{Z}/2^k\mathbb{Z}$ projiziert, ist ein
projektives System. Ihr projektiver Limes $\pi_\infty$ ist ein $T$-invariantes Maß auf
$\mathbb{Z}_2$.
\end{theorem}

\begin{proof}
Dies folgt aus der Kompaktheit von $\mathbb{Z}_2$ und der Konsistenz der Projektionen.
Der projektive Limes einer konsistenten Folge von Wahrscheinlichkeitsmaßen auf
endlichen Mengen ist ein Wahrscheinlichkeitsmaß auf dem projektiven Limes $\mathbb{Z}_2$.
Die $T$-Invarianz wird von der Invarianz jedes $\pi_k$ unter der mod-$2^k$-Collatz-Abbildung
geerbt.
\end{proof}

\section{Der ergodische Ansatz für nicht-triviale Zyklen}

\begin{theorem}[Fehlen nicht-trivialer Zyklen aus ergodischer Sicht]\label{thm:keine_zyklen}
Ein nicht-trivialer Zyklus von $T$ in $\mathbb{N}$ (falls er existierte) würde einer
periodischen Bahn der Collatz-Abbildung auf einer Maß-Null-Teilmenge von $\mathbb{Z}_2$
entsprechen. Seine Existenz würde nicht den ergodischen Eigenschaften von $T$ auf
$(\mathbb{Z}_2, \mu)$ widersprechen, da einzelne Bahnen sich anders verhalten können
als $\mu$-fast-jede Bahn.
\end{theorem}

\begin{proof}
Ein Zyklus $\{n_0, n_1, \ldots, n_{k-1}\}$ in $\mathbb{N}$ hat das Haar-Maß null in
$\mathbb{Z}_2$ (als endliche Menge). Der Ergodensatz macht Aussagen über $\mu$-fast-jede
Bahn, was individuelle Bahnen nicht einschränkt. Daher muss die Nicht-Existenz von
Zyklen durch arithmetische Methoden (vgl.\ die Zyklusgleichung aus Paper~29) nachgewiesen
werden, nicht durch ergodische Methoden.
\end{proof}

\begin{remark}[Was Ergodentheorie kann und nicht kann]\label{rem:ergodisch_grenzen}
Ergodentheorie liefert:
\begin{itemize}
  \item Evidenz, dass ``generische'' (maßtypische) Bahnen konvergieren.
  \item Den negativen Drift $\log(3/4)$ pro Syrakus-Schritt.
  \item Mischung und Abfall von Korrelationen für die 2-adische Dynamik.
\end{itemize}
Aber sie kann nicht beweisen:
\begin{itemize}
  \item Dass jede \emph{spezifische} Bahn konvergiert (da positive Zahlen Maß null haben).
  \item Das Fehlen nicht-trivialer Zyklen (ein Maß-Null-Phänomen).
  \item Die vollständige Collatz-Vermutung ohne zusätzlichen arithmetischen Input.
\end{itemize}
\end{remark}

\section{Zusammenfassung}

Wir haben die ergodisch-theoretischen Grundlagen des Collatz-Problems entwickelt:
\begin{enumerate}
  \item Das Collatz-dynamische System auf $\mathbb{Z}_2$ ist nicht maßerhaltend,
        lässt aber invariante Maße zu.
  \item Die Syrakus-Abbildung $S$ ist ergodisch und mischend mit exponentiellem Abfall
        von Korrelationen.
  \item Birkhoffs Ergodensatz liefert fast sichere Konvergenz zu $\log(3/4) < 0$ pro
        Schritt und bestätigt den negativen Drift.
  \item Furstenbergs Prinzip verbindet Dichteresultate mit ergodischen Eigenschaften.
  \item Die Verbindung zu Bedford--McMullen-Maßen und symbolischer Dynamik bietet einen
        geometrischen Rahmen für den Raum der Collatz-Bahntypen.
  \item Stochastische Matrizen (Collatz-Markov-Ketten) liefern eine endlichdimensionale
        Approximation der vollen Dynamik.
\end{enumerate}
Die ergodische Sichtweise legt die Wahrheit der Collatz-Vermutung nahe, kann sie aber
nicht beweisen; zusätzliche arithmetische und $p$-adische Methoden sind erforderlich.

\begin{thebibliography}{99}

\bibitem{Collatz1937}
L.~Collatz.
\newblock Unveröffentlichtes Problem, vorgestellt beim Internationalen Mathematikerkongress, 1937.

\bibitem{Lagarias1985}
J.~C. Lagarias.
\newblock The $3x+1$ problem and its generalizations.
\newblock \emph{American Mathematical Monthly}, 92(1):3--23, 1985.

\bibitem{Lagarias2010}
J.~C. Lagarias (Hrsg.).
\newblock \emph{The Ultimate Challenge: The $3x+1$ Problem}.
\newblock American Mathematical Society, Providence, RI, 2010.

\bibitem{Tao2022}
T.~Tao.
\newblock Almost all Collatz orbits attain almost bounded values.
\newblock \emph{Forum of Mathematics, Sigma}, 10:e72, 2022.

\bibitem{Furstenberg1977}
H.~Furstenberg.
\newblock Ergodic behavior of diagonal measures and a theorem of Szemer\'edi on arithmetic progressions.
\newblock \emph{Journal d'Analyse Math\'ematique}, 31(1):204--256, 1977.

\bibitem{Walters1982}
P.~Walters.
\newblock \emph{An Introduction to Ergodic Theory}.
\newblock Springer-Verlag, New York, 1982.

\bibitem{Bedford1984}
T.~Bedford.
\newblock Crinkly curves, Markov partitions and box dimensions in self-similar sets.
\newblock Dissertation, University of Warwick, 1984.

\bibitem{McMullen1984}
C.~McMullen.
\newblock The Hausdorff dimension of general Sierpi\'nski carpets.
\newblock \emph{Nagoya Mathematical Journal}, 96:1--9, 1984.

\bibitem{Koblitz1984}
N.~Koblitz.
\newblock \emph{$p$-adic Numbers, $p$-adic Analysis, and Zeta Functions}.
\newblock Springer-Verlag, New York, 1984.

\end{thebibliography}

\end{document}
